% GENERATED FILE. Edit canonical lessons, then run npm run generate:books.
% canonical-source-hash: fnv1a64:e24991c23c1d8b63
% canonical-lessons: JA-C266-oboreru, JA-C266-tobikomu, JA-C266-hipparu, JA-C266-yatto, JA-C266-buji

\chapter{Trouble in the River}
\label{ch:ja-trouble-in-the-river}
\hlchaptermodality{266}

\begin{chapteropening}
\textbf{By the end of this chapter:} \emph{I can say who fell in the water and who helped.}

Complete the last lesson of chapter 266: I can say who fell in the water and who helped.
\end{chapteropening}

\section[\texorpdfstring{\ja{おぼれる} (oboreru)}{おぼれる (oboreru)}]{\ja{おぼれる} (oboreru) — to drown; to be drowning}

\label{lesson:JA-C266-oboreru}



\begin{quote}
Before the new one: say the Japanese for an escape route, then the Japanese for to rush to the scene.
\end{quote}

\subsection*{You'll want to know: \ja{おぼれる}}
\textbf{\ja{おぼれる}} — \emph{oboreru} — \textquotedblleft{}to drown; to be drowning\textquotedblright{}.

\begin{quote}
\textbf{\ja{こどもが} \ja{うみで} \ja{おぼれました。}} — \emph{kodomo ga umi de oboremashita} — \textquotedblleft{}A child got into trouble in the sea.\textquotedblright{}
\end{quote}

Polite \textbf{\ja{おぼれます}}, \emph{oboremasu}; an ichidan verb. Past (polite): \textbf{\ja{おとこが} \ja{おぼれました。}}, \emph{otoko ga oboremashita}.

\begin{cousinweb}
\textbf{\ja{おぼれる}} (\emph{oboreru}) is a native Japanese word (wago), inherited rather than borrowed from Chinese; its kanji was attached later as a kun'yomi (JMdict).
\end{cousinweb}

\begin{grammarlens}[title={What you've built}]
One more word for telling what happened.
\end{grammarlens}

\subsection*{Your turn}
\begin{itemize}
\raggedright
  \item \emph{Say it:} \emph{oboreru}
  \item \emph{Say it:} \emph{oboreru}, once more
  \item \emph{Say it:} \emph{kodomo ga umi de oboremashita} — A child got into trouble in the sea.
  \item \emph{Recall:} say the Japanese for an escape route, then the Japanese for to rush to the scene, then say \emph{oboreru} again
\end{itemize}

\begin{tcolorbox}[breakable,colback=teal!4,colframe=teal!35!black,title={Before you move on}]
What does \ja{おぼれる} mean? (\textquotedblleft{}To drown; to be drowning\textquotedblright{}.) Say it once more.
\end{tcolorbox}
\section[\texorpdfstring{\ja{とびこむ} (tobikomu)}{とびこむ (tobikomu)}]{\ja{とびこむ} (tobikomu) — to jump in; to dive in}

\label{lesson:JA-C266-tobikomu}



\begin{quote}
Before the new one: say the Japanese for to rush to the scene, then the Japanese for to drown.
\end{quote}

\subsection*{You'll want to know: \ja{とびこむ}}
\textbf{\ja{とびこむ}} — \emph{tobikomu} — \textquotedblleft{}to jump in; to dive in\textquotedblright{}.

\begin{quote}
\textbf{\ja{こどもが} \ja{おぼれたら、すぐ} \ja{とびこみます。}} — \emph{kodomo ga oboretara, sugu tobikomimasu} — \textquotedblleft{}If a child is drowning, I jump in at once.\textquotedblright{}
\end{quote}

Polite \textbf{\ja{とびこみます}}, \emph{tobikomimasu}; a godan verb. Past (polite): \textbf{\ja{ちちが} \ja{かわに} \ja{とびこみました。}}, \emph{chichi ga kawa ni tobikomimashita}.

\begin{cousinweb}
\textbf{\ja{とびこむ}} (\emph{tobikomu}) is a native Japanese word (wago), inherited rather than borrowed from Chinese, built from \emph{tobi}, \textquotedblleft{}leaping\textquotedblright{} (from \emph{tobu}), and \emph{komu}, \textquotedblleft{}to go into\textquotedblright{} (JMdict, KANJIDIC2).
\end{cousinweb}

\begin{grammarlens}[title={What you've built}]
One more word for telling what happened.
\end{grammarlens}

\subsection*{Your turn}
\begin{itemize}
\raggedright
  \item \emph{Say it:} \emph{tobikomu}
  \item \emph{Say it:} \emph{tobikomu}, once more
  \item \emph{Say it:} \emph{kodomo ga oboretara, sugu tobikomimasu} — If a child is drowning, I jump in at once.
  \item \emph{Recall:} say the Japanese for to rush to the scene, then the Japanese for to drown, then say \emph{tobikomu} again
\end{itemize}

\begin{tcolorbox}[breakable,colback=teal!4,colframe=teal!35!black,title={Before you move on}]
What does \ja{とびこむ} mean? (\textquotedblleft{}To jump in; to dive in\textquotedblright{}.) Say it once more.
\end{tcolorbox}
\section[\texorpdfstring{\ja{ひっぱる} (hipparu)}{ひっぱる (hipparu)}]{\ja{ひっぱる} (hipparu) — to tug; to pull hard}

\label{lesson:JA-C266-hipparu}



\begin{quote}
Before the new one: say the Japanese for to drown, then the Japanese for to jump in.
\end{quote}

\subsection*{You'll want to know: \ja{ひっぱる}}
\textbf{\ja{ひっぱる}} — \emph{hipparu} — \textquotedblleft{}to tug; to pull hard\textquotedblright{}.

\begin{quote}
\textbf{\ja{ちちが} \ja{こどもの} \ja{てを} \ja{ひっぱりました。}} — \emph{chichi ga kodomo no te o hipparimashita} — \textquotedblleft{}My father pulled the child by the hand.\textquotedblright{}
\end{quote}

Polite \textbf{\ja{ひっぱります}}, \emph{hipparimasu}; a godan verb. Past (polite): \textbf{\ja{あにが} \ja{なわを} \ja{ひっぱりました。}}, \emph{ani ga nawa o hipparimashita}.

\begin{cousinweb}
\textbf{\ja{ひっぱる}} (\emph{hipparu}) is a native Japanese word (wago), inherited rather than borrowed from Chinese, built from \emph{hik-}, \textquotedblleft{}pulling\textquotedblright{} (from \emph{hiku}), and \emph{haru}, \textquotedblleft{}to stretch\textquotedblright{} (JMdict, KANJIDIC2).
\end{cousinweb}

\begin{grammarlens}[title={What you've built}]
One more word for telling what happened.
\end{grammarlens}

\subsection*{Your turn}
\begin{itemize}
\raggedright
  \item \emph{Say it:} \emph{hipparu}
  \item \emph{Say it:} \emph{hipparu}, once more
  \item \emph{Say it:} \emph{chichi ga kodomo no te o hipparimashita} — My father pulled the child by the hand.
  \item \emph{Recall:} say the Japanese for to drown, then the Japanese for to jump in, then say \emph{hipparu} again
\end{itemize}

\begin{tcolorbox}[breakable,colback=teal!4,colframe=teal!35!black,title={Before you move on}]
What does \ja{ひっぱる} mean? (\textquotedblleft{}To tug; to pull hard\textquotedblright{}.) Say it once more.
\end{tcolorbox}
\section[\texorpdfstring{\ja{やっと} (yatto)}{やっと (yatto)}]{\ja{やっと} (yatto) — at last; finally}

\label{lesson:JA-C266-yatto}



\begin{quote}
Before the new one: say the Japanese for to jump in, then the Japanese for to tug.
\end{quote}

\subsection*{You'll want to know: \ja{やっと}}
\textbf{\ja{やっと}} — \emph{yatto} — \textquotedblleft{}at last; finally\textquotedblright{}.

\begin{quote}
\textbf{\ja{こどもは} \ja{やっと} \ja{たすかりました。}} — \emph{kodomo wa yatto tasukarimashita} — \textquotedblleft{}In the end the child was saved.\textquotedblright{}
\end{quote}

\begin{cousinweb}
\textbf{\ja{やっと}} (\emph{yatto}) is a native Japanese mimetic word, inherited rather than borrowed from Chinese, and its older origin is not settled (JMdict).
\end{cousinweb}

\begin{grammarlens}[title={What you've built}]
One more word for telling what happened.
\end{grammarlens}

\subsection*{Your turn}
\begin{itemize}
\raggedright
  \item \emph{Say it:} \emph{yatto}
  \item \emph{Say it:} \emph{yatto}, once more
  \item \emph{Say it:} \emph{kodomo wa yatto tasukarimashita} — In the end the child was saved.
  \item \emph{Recall:} say the Japanese for to jump in, then the Japanese for to tug, then say \emph{yatto} again
\end{itemize}

\begin{tcolorbox}[breakable,colback=teal!4,colframe=teal!35!black,title={Before you move on}]
What does \ja{やっと} mean? (\textquotedblleft{}At last; finally\textquotedblright{}.) Say it once more.
\end{tcolorbox}
\section[\texorpdfstring{\ja{ぶじ} (buji)}{ぶじ (buji)}]{\ja{ぶじ} (buji) — unharmed; safe and sound}

\label{lesson:JA-C266-buji}



\begin{quote}
Before the new one: say the Japanese for to tug, then the Japanese for at last.
\end{quote}

\subsection*{You'll want to know: \ja{ぶじ}}
\textbf{\ja{ぶじ}} — \emph{buji} — \textquotedblleft{}unharmed; safe and sound\textquotedblright{}.

\begin{quote}
\textbf{\ja{こどもは} \ja{ぶじでした。}} — \emph{kodomo wa buji deshita} — \textquotedblleft{}The child was unhurt.\textquotedblright{}
\end{quote}

A na-adjective: \ja{な} before a noun.

\begin{cousinweb}
\textbf{\ja{ぶじ}} (\emph{buji}) is Sino-Japanese (kango): two Chinese characters read with their on'yomi, readings borrowed from Chinese. They are \emph{bu}, \textquotedblleft{}nothing\textquotedblright{} (Middle Chinese \emph{mju}) and \emph{ji}, \textquotedblleft{}matter\textquotedblright{} (\emph{dzriH}) (KANJIDIC2; Middle Chinese in Baxter's transcription).
\end{cousinweb}

\begin{grammarlens}[title={What you've built}]
One more word for telling what happened.
\end{grammarlens}

\subsection*{Your turn}
\begin{itemize}
\raggedright
  \item \emph{Say it:} \emph{buji}
  \item \emph{Say it:} \emph{buji}, once more
  \item \emph{Say it:} \emph{kodomo wa buji deshita} — The child was unhurt.
  \item \emph{Recall:} say the Japanese for to tug, then the Japanese for at last, then say \emph{buji} again
\end{itemize}

\begin{tcolorbox}[breakable,colback=teal!4,colframe=teal!35!black,title={Before you move on}]
What does \ja{ぶじ} mean? (\textquotedblleft{}Unharmed; safe and sound\textquotedblright{}.) Say it once more.
\end{tcolorbox}
